% !TEX root = ../main.tex
\section{Experimental setup and main results}
\label{sec:experiments}
\paragraph{Tasks and models.}
We evaluate on the 1,319-problem GSM8K test set \citep{gsm8k} and the 500-problem MATH500 benchmark drawn from MATH \citep{math,math500}. GSM8K training uses 2.43 million solution rows combining GSM-style subsets of OpenMathInstruct-2, Orca-Math, and MetaMathQA \citep{openmath,orcamath,metamath}; the MATH model uses 3.97 million eligible MATH-family solution rows from OpenMathInstruct-2. All models use a 1,024-token canvas and the native Qwen tokenizer. The problem is supplied without demonstrations, followed by an instruction to reason step by step and put the final answer in a box. Final learned-prompt inference does not execute Qwen Transformer layers.

ELF-B has 12 blocks of width 768; ELF-L has 32 blocks of width 1,280. Their denoising paths use 90.44M and 638.15M parameters, respectively. Terminal-only decoder parameters add 156.52M and 157.05M, bringing the shared denoiser/decoder modules to 246.96M and 795.21M. The prompt encoders add 45.02M and 121.33M parameters, and both retain a frozen 388.96M-element Qwen lookup. Prompt encoding occurs once; the shared backbone runs at every denoising step and once more for terminal decoding. Appendix~\ref{app:implementation} gives the complete accounting.

\paragraph{Training and endpoint selection.}
We train the denoiser with teacher prompt conditioning for 12 epochs, then jointly adapt the learned prompt encoder and denoiser through epoch 18 for ELF-B and epoch 21 for ELF-L. The effective supervised batch is 512, with constant learning rate $0.002$, Muon and auxiliary optimizers, zero weight decay, and gradient clipping at one. We select NFT updates 300, 500, and 600 for GSM8K ELF-B, GSM8K ELF-L, and MATH ELF-L, respectively, using fixed development sets of 1,000 GSM-style and 500 MATH problems. Selection uses two-seed mean accuracy and favors the earlier update on an exact tie.

\paragraph{Evaluation.}
We first evaluate four sampling seeds at 32 and 64 steps. Each supervised/NFT pair uses a common step budget chosen by its equally weighted pilot mean pass@1; ties favor 32. All three pairs select 64 steps. Four further seeds give eight samples per problem at that budget. The stage and schedule ablations use 32 steps throughout. Supervised evaluation selects ELF/prompt EMA $0.9999/0.9999$; NFT selects $0.99/0.99$. The immediate prompt-swap experiment selects the standalone MSE prompt EMA $0.999$. GSM8K uses canonical numerical-answer scoring; MATH uses Math-Verify, with PRM800K retained as a secondary scorer \citep{mathverify,prm800k}.

We report mean pass@1 and the standard unbiased pass@$k$ estimator \citep{codex}. Confidence intervals use 4,000 paired problem-bootstrap replicates, retaining matched problems across conditions. They quantify uncertainty for fixed trained checkpoints and the observed sampling seeds. Appendix~\ref{app:evaluation} specifies the estimators, seed lists, and selection procedure.

\subsection{Reasoning performance}
% !TEX root = ../main.tex
\begin{table}[t]
\caption{Main results at 64 denoising steps with eight samples per problem. Accuracy is in percent; brackets are 95\% problem-bootstrap intervals for mean pass@1. All rows use clocks $(2.5,2,1.5)$ and SCCFG 2. The matched MATH500 SCCFG 3 comparison appears in Appendix~\ref{app:mathsccfg3}.}
\label{tab:main}
\centering\small
\begin{tabular}{llrrrr}
\toprule
Benchmark / model & Stage & Pass@1 [95\% CI] & Pass@2 & Pass@4 & Pass@8 \\
\midrule
GSM8K / ELF-B & Pre-NFT & 39.17 [37.17, 41.20] & 49.82 & 59.20 & 67.32 \\
GSM8K / ELF-B & Post-NFT & 42.16 [39.97, 44.33] & 51.36 & 59.40 & 66.26 \\
GSM8K / ELF-L & Pre-NFT & 58.67 [56.45, 60.78] & 68.17 & 75.13 & 80.21 \\
GSM8K / ELF-L & Post-NFT & 63.74 [61.52, 65.88] & 71.48 & 77.34 & 81.96 \\
MATH500 / ELF-L & Pre-NFT & 18.85 [16.42, 21.40] & 27.50 & 37.13 & 47.20 \\
MATH500 / ELF-L & Post-NFT & 23.60 [20.80, 26.45] & 33.00 & 43.07 & 52.60 \\
\bottomrule
\end{tabular}
\end{table}

Table~\ref{tab:main} reports the six main endpoints. NFT raises mean pass@1 from 39.17\% to 42.16\% for GSM8K ELF-B, from 58.67\% to 63.74\% for GSM8K ELF-L, and from 18.85\% to 23.60\% for MATH ELF-L. The effect on oracle coverage is less uniform. At eight samples, GSM8K ELF-B changes from 67.32\% to 66.26\%, whereas MATH increases from 47.20\% to 52.60\%. Thus the benefit of post-training is not summarized fully by a single-sample score.

\begin{table}[t]
\caption{Reported single-sample GSM8K results from recent continuous language models. External rows retain their papers' training and evaluation protocols; they are contextual comparisons rather than matched-data ablations. Our rows use eight-seed mean pass@1 at 64 steps after NFT.}
\label{tab:literature}
\centering\small
\begin{tabular}{llr}
\toprule
Method & Reported setting & Accuracy (\%)\\
\midrule
$\mathbb S$-FLM \citep{hyperspherical} & Top-1 velocity, $T=1$, 1024 steps & 18.00\\
FMLM+ (Init) \citep{posterior} & Posterior refinement, 64 NFEs & 33.60\\
MLFM \citep{mlfm} & Guidance + online token promotion & 31.24\\
\midrule
ELF-B + NFT & Natural-language solutions, 64 steps & 42.16\\
ELF-L + NFT & Natural-language solutions, 64 steps & 63.74\\
\bottomrule
\end{tabular}
\end{table}

Table~\ref{tab:literature} compares continuous models using the original papers' reported accuracies. Our GSM8K results exceed the listed operating points, with different supervision and decoding protocols. This advantage is present below our headline budget: ELF-B reaches 37.64\% at 16 denoising steps (17 backbone calls, seeds 42/123), versus FMLM+ Init's reported 33.6\% at 64 NFEs. The complete depth comparison in Appendix~\ref{app:literaturecontinuous} retains both low- and high-budget results; our current sweep does not establish performance below eight denoising steps. No MATH500 result is reported for these three continuous comparators. Discrete masked and blockwise models form a separate comparison group in Appendix~\ref{app:literaturediscrete}, alongside their sizes and published protocols.

\subsection{Representation controls}
\label{sec:representationcontrols}
We separate projector construction from layer selection using three six-epoch ELF-B arms: learned projectors over layers 16/24/32, PCA over the same layers and dimensions, and a 1,024-dimensional PCA representation from layer 32 alone. All representations are whitened. The first pair fixes layer selection; the second fixes the projector family and total latent width. Their layout is shown in Table~\ref{tab:representation}. Epoch-six evaluation uses a common EMA, seed set, and inference schedule.

\begin{table}[t]
\caption{Representation ablation at six GSM8K epochs. The learned reference uses original Qwen prompt conditioning. All entries use ELF EMA $0.999$ and synchronous 32-step inference with seeds 42/123. Dashes reserve the pending PCA results.}
\label{tab:representation}
\centering\small
\begin{tabular}{llcr}
\toprule
Projector & Teacher layers & Latent widths & Pass@1 (\%)\\
\midrule
Learned, whitened & 16/24/32 & 256/256/512 & 26.76\\
PCA, whitened & 16/24/32 & 256/256/512 & \pending\\
PCA, whitened & 32 & 1024 & \pending\\
\bottomrule
\end{tabular}
\end{table}
